\subsection{コーシーの平均値の定理}
素朴に考えれば,\,$f$ と $g$ のそれぞれに平均値の定理を使い,2つの式の比をとればよさそうである.しかし,そうして得られるのは $\frac{f'(c_1)}{g'(c_2)}$ という比であり,\,$f$ に対する点 $c_1$ と $g$ に対する点 $c_2$ は一般には異なる.実は,2つの導関数を同じ点で比べた形の等式が成り立つ.以下では,区間の端を $u,v$,区間の内側の点を $\xi$ (グザイ)と書く.

\begin{theorem}{コーシーの平均値の定理}\label{定理:コーシー平均値}
$u<v$ とする.\,$f,g$ が閉区間 $[u,v]$ で連続,開区間 $(u,v)$ で微分可能で,\,$(u,v)$ のすべての点で $g'\ne0$ なら,
\[
 \frac{f(v)-f(u)}{g(v)-g(u)}=\frac{f'(\xi)}{g'(\xi)}
\]
を満たす $\xi\in(u,v)$ が少なくとも1つ存在する.
\end{theorem}

$g(x)=x$ とすれば $g(v)-g(u)=v-u$,\,$g'=1$ だから,これは平均値の定理そのものである.

この定理も,平均値の定理と同じく,ロルの定理(定理\ref{定理:ロル})から導ける.まず,左辺の分母が $0$ でないことを確かめておこう.もし $g(v)=g(u)$ であれば,ロルの定理から $(u,v)$ のどこかで $g'=0$ となってしまい,仮定に反する.だから $g(v)-g(u)\ne0$ である.

そこで,両端で値が等しくなる関数
\[
 H(t)=\{f(v)-f(u)\}\{g(t)-g(u)\}
      -\{g(v)-g(u)\}\{f(t)-f(u)\}
\]
を作る.\,$t=u$ を代入すると,2つの項はどちらも $0$ になる.\,$t=v$ を代入すると,2つの項は同じ積になって打ち消し合う.すなわち
\[
 H(u)=H(v)=0
\]
である.\,$H$ は $[u,v]$ で連続,\,$(u,v)$ で微分可能だから,ロルの定理により $H'(\xi)=0$ となる点 $\xi\in(u,v)$ がある.\,$H$ を $t$ で微分すると,\,$f(u),f(v),g(u),g(v)$ は定数なので,
\[
 H'(\xi)=\{f(v)-f(u)\}g'(\xi)-\{g(v)-g(u)\}f'(\xi)=0
\]
となる.\,$g(v)-g(u)$ も $g'(\xi)$ も $0$ でないから,両辺を $\{g(v)-g(u)\}g'(\xi)$ で割って移項すれば,求める等式を得る.
